\IfFileExists{stacks-project.cls}{%
\documentclass{stacks-project}
}{%
\documentclass{amsart}
}

% For dealing with references we use the comment environment
\usepackage{verbatim}
\newenvironment{reference}{\comment}{\endcomment}
%\newenvironment{reference}{}{}
\newenvironment{slogan}{\comment}{\endcomment}
\newenvironment{history}{\comment}{\endcomment}

% For commutative diagrams we use Xy-pic
\usepackage[all]{xy}

% We use 2cell for 2-commutative diagrams.
\xyoption{2cell}
\UseAllTwocells

% We use multicol for the list of chapters between chapters
\usepackage{multicol}

% This is generally recommended for better output
\usepackage{lmodern}
\usepackage[T1]{fontenc}

% For cross-file-references
\usepackage{xr-hyper}

% Package for hypertext links:
\usepackage{hyperref}

% For any local file, say "hello.tex" you want to link to please
% use \externaldocument[hello-]{hello}
\externaldocument[introduction-]{introduction}
\externaldocument[conventions-]{conventions}
\externaldocument[sets-]{sets}
\externaldocument[categories-]{categories}
\externaldocument[topology-]{topology}
\externaldocument[sheaves-]{sheaves}
\externaldocument[sites-]{sites}
\externaldocument[stacks-]{stacks}
\externaldocument[fields-]{fields}
\externaldocument[algebra-]{algebra}
\externaldocument[brauer-]{brauer}
\externaldocument[homology-]{homology}
\externaldocument[derived-]{derived}
\externaldocument[simplicial-]{simplicial}
\externaldocument[more-algebra-]{more-algebra}
\externaldocument[smoothing-]{smoothing}
\externaldocument[modules-]{modules}
\externaldocument[sites-modules-]{sites-modules}
\externaldocument[injectives-]{injectives}
\externaldocument[cohomology-]{cohomology}
\externaldocument[sites-cohomology-]{sites-cohomology}
\externaldocument[dga-]{dga}
\externaldocument[dpa-]{dpa}
\externaldocument[sdga-]{sdga}
\externaldocument[hypercovering-]{hypercovering}
\externaldocument[schemes-]{schemes}
\externaldocument[constructions-]{constructions}
\externaldocument[properties-]{properties}
\externaldocument[morphisms-]{morphisms}
\externaldocument[coherent-]{coherent}
\externaldocument[divisors-]{divisors}
\externaldocument[limits-]{limits}
\externaldocument[varieties-]{varieties}
\externaldocument[topologies-]{topologies}
\externaldocument[descent-]{descent}
\externaldocument[perfect-]{perfect}
\externaldocument[more-morphisms-]{more-morphisms}
\externaldocument[flat-]{flat}
\externaldocument[groupoids-]{groupoids}
\externaldocument[more-groupoids-]{more-groupoids}
\externaldocument[etale-]{etale}
\externaldocument[chow-]{chow}
\externaldocument[intersection-]{intersection}
\externaldocument[pic-]{pic}
\externaldocument[weil-]{weil}
\externaldocument[adequate-]{adequate}
\externaldocument[dualizing-]{dualizing}
\externaldocument[duality-]{duality}
\externaldocument[discriminant-]{discriminant}
\externaldocument[derham-]{derham}
\externaldocument[local-cohomology-]{local-cohomology}
\externaldocument[algebraization-]{algebraization}
\externaldocument[curves-]{curves}
\externaldocument[resolve-]{resolve}
\externaldocument[models-]{models}
\externaldocument[functors-]{functors}
\externaldocument[equiv-]{equiv}
\externaldocument[pione-]{pione}
\externaldocument[etale-cohomology-]{etale-cohomology}
\externaldocument[proetale-]{proetale}
\externaldocument[relative-cycles-]{relative-cycles}
\externaldocument[more-etale-]{more-etale}
\externaldocument[trace-]{trace}
\externaldocument[crystalline-]{crystalline}
\externaldocument[spaces-]{spaces}
\externaldocument[spaces-properties-]{spaces-properties}
\externaldocument[spaces-morphisms-]{spaces-morphisms}
\externaldocument[decent-spaces-]{decent-spaces}
\externaldocument[spaces-cohomology-]{spaces-cohomology}
\externaldocument[spaces-limits-]{spaces-limits}
\externaldocument[spaces-divisors-]{spaces-divisors}
\externaldocument[spaces-over-fields-]{spaces-over-fields}
\externaldocument[spaces-topologies-]{spaces-topologies}
\externaldocument[spaces-descent-]{spaces-descent}
\externaldocument[spaces-perfect-]{spaces-perfect}
\externaldocument[spaces-more-morphisms-]{spaces-more-morphisms}
\externaldocument[spaces-flat-]{spaces-flat}
\externaldocument[spaces-groupoids-]{spaces-groupoids}
\externaldocument[spaces-more-groupoids-]{spaces-more-groupoids}
\externaldocument[bootstrap-]{bootstrap}
\externaldocument[spaces-pushouts-]{spaces-pushouts}
\externaldocument[spaces-chow-]{spaces-chow}
\externaldocument[groupoids-quotients-]{groupoids-quotients}
\externaldocument[spaces-more-cohomology-]{spaces-more-cohomology}
\externaldocument[spaces-simplicial-]{spaces-simplicial}
\externaldocument[spaces-duality-]{spaces-duality}
\externaldocument[formal-spaces-]{formal-spaces}
\externaldocument[restricted-]{restricted}
\externaldocument[spaces-resolve-]{spaces-resolve}
\externaldocument[formal-defos-]{formal-defos}
\externaldocument[defos-]{defos}
\externaldocument[cotangent-]{cotangent}
\externaldocument[examples-defos-]{examples-defos}
\externaldocument[algebraic-]{algebraic}
\externaldocument[examples-stacks-]{examples-stacks}
\externaldocument[stacks-sheaves-]{stacks-sheaves}
\externaldocument[criteria-]{criteria}
\externaldocument[artin-]{artin}
\externaldocument[quot-]{quot}
\externaldocument[stacks-properties-]{stacks-properties}
\externaldocument[stacks-morphisms-]{stacks-morphisms}
\externaldocument[stacks-limits-]{stacks-limits}
\externaldocument[stacks-cohomology-]{stacks-cohomology}
\externaldocument[stacks-perfect-]{stacks-perfect}
\externaldocument[stacks-introduction-]{stacks-introduction}
\externaldocument[stacks-more-morphisms-]{stacks-more-morphisms}
\externaldocument[stacks-geometry-]{stacks-geometry}
\externaldocument[moduli-]{moduli}
\externaldocument[moduli-curves-]{moduli-curves}
\externaldocument[examples-]{examples}
\externaldocument[exercises-]{exercises}
\externaldocument[guide-]{guide}
\externaldocument[desirables-]{desirables}
\externaldocument[coding-]{coding}
\externaldocument[obsolete-]{obsolete}
\externaldocument[fdl-]{fdl}
\externaldocument[index-]{index}

% Theorem environments.
%
\theoremstyle{plain}
\newtheorem{theorem}[subsection]{Theorem}
\newtheorem{proposition}[subsection]{Proposition}
\newtheorem{lemma}[subsection]{Lemma}

\theoremstyle{definition}
\newtheorem{definition}[subsection]{Definition}
\newtheorem{example}[subsection]{Example}
\newtheorem{exercise}[subsection]{Exercise}
\newtheorem{situation}[subsection]{Situation}

\theoremstyle{remark}
\newtheorem{remark}[subsection]{Remark}
\newtheorem{remarks}[subsection]{Remarks}

\numberwithin{equation}{subsection}

% Macros
%
\def\lim{\mathop{\mathrm{lim}}\nolimits}
\def\colim{\mathop{\mathrm{colim}}\nolimits}
\def\Spec{\mathop{\mathrm{Spec}}}
\def\Hom{\mathop{\mathrm{Hom}}\nolimits}
\def\Ext{\mathop{\mathrm{Ext}}\nolimits}
\def\SheafHom{\mathop{\mathcal{H}\!\mathit{om}}\nolimits}
\def\SheafExt{\mathop{\mathcal{E}\!\mathit{xt}}\nolimits}
\def\Sch{\mathit{Sch}}
\def\Mor{\mathop{\mathrm{Mor}}\nolimits}
\def\Ob{\mathop{\mathrm{Ob}}\nolimits}
\def\Sh{\mathop{\mathit{Sh}}\nolimits}
\def\NL{\mathop{N\!L}\nolimits}
\def\CH{\mathop{\mathrm{CH}}\nolimits}
\def\proetale{{pro\text{-}\acute{e}tale}}
\def\etale{{\acute{e}tale}}
\def\QCoh{\mathit{QCoh}}
\def\Ker{\mathop{\mathrm{Ker}}}
\def\Im{\mathop{\mathrm{Im}}}
\def\Coker{\mathop{\mathrm{Coker}}}
\def\Coim{\mathop{\mathrm{Coim}}}

% Boxtimes
%
\DeclareMathSymbol{\boxtimes}{\mathbin}{AMSa}{"02}

%
% Macros for moduli stacks/spaces
%
\def\QCohstack{\mathcal{QC}\!\mathit{oh}}
\def\Cohstack{\mathcal{C}\!\mathit{oh}}
\def\Spacesstack{\mathcal{S}\!\mathit{paces}}
\def\Quotfunctor{\mathrm{Quot}}
\def\Hilbfunctor{\mathrm{Hilb}}
\def\Curvesstack{\mathcal{C}\!\mathit{urves}}
\def\Polarizedstack{\mathcal{P}\!\mathit{olarized}}
\def\Complexesstack{\mathcal{C}\!\mathit{omplexes}}
% \Pic is the operator that assigns to X its picard group, usage \Pic(X)
% \Picardstack_{X/B} denotes the Picard stack of X over B
% \Picardfunctor_{X/B} denotes the Picard functor of X over B
\def\Pic{\mathop{\mathrm{Pic}}\nolimits}
\def\Picardstack{\mathcal{P}\!\mathit{ic}}
\def\Picardfunctor{\mathrm{Pic}}
\def\Deformationcategory{\mathcal{D}\!\mathit{ef}}
\usepackage{cleveref}
\crefname{originaldefinition}{original definition}{original definitions}
\Crefname{originaldefinition}{Original definition}{Original definitions}
\makeatletter
\newlabel{descent-definition-descent-datum}{{023V}{source}{}{sourcepointer-1}{}}
\newlabel{descent-definition-descent-datum@cref}{{[originaldefinition][0][]023V}{[1][0][]source}{}{}{}}
\makeatother
\crefname{originallemma}{original lemma}{original lemmas}
\Crefname{originallemma}{Original lemma}{Original lemmas}
\makeatletter
\newlabel{descent-lemma-descending-property-etale}{{02VN}{source}{}{sourcepointer-2}{}}
\newlabel{descent-lemma-descending-property-etale@cref}{{[originallemma][0][]02VN}{[1][0][]source}{}{}{}}
\makeatother
\crefname{originallemma}{original lemma}{original lemmas}
\Crefname{originallemma}{Original lemma}{Original lemmas}
\makeatletter
\newlabel{descent-lemma-descending-property-isomorphism}{{02L4}{source}{}{sourcepointer-3}{}}
\newlabel{descent-lemma-descending-property-isomorphism@cref}{{[originallemma][0][]02L4}{[1][0][]source}{}{}{}}
\makeatother
\crefname{originallemma}{original lemma}{original lemmas}
\Crefname{originallemma}{Original lemma}{Original lemmas}
\makeatletter
\newlabel{descent-lemma-descending-property-surjective}{{02KV}{source}{}{sourcepointer-4}{}}
\newlabel{descent-lemma-descending-property-surjective@cref}{{[originallemma][0][]02KV}{[1][0][]source}{}{}{}}
\makeatother
\crefname{originallemma}{original lemma}{original lemmas}
\Crefname{originallemma}{Original lemma}{Original lemmas}
\makeatletter
\newlabel{descent-lemma-fpqc-universal-effective-epimorphisms}{{023Q}{source}{}{sourcepointer-5}{}}
\newlabel{descent-lemma-fpqc-universal-effective-epimorphisms@cref}{{[originallemma][0][]023Q}{[1][0][]source}{}{}{}}
\makeatother
\crefname{originallemma}{original lemma}{original lemmas}
\Crefname{originallemma}{Original lemma}{Original lemmas}
\makeatletter
\newlabel{lemma-diagonal-formal-algebraic-space}{{0AIP}{source}{}{sourcepointer-6}{}}
\newlabel{lemma-diagonal-formal-algebraic-space@cref}{{[originallemma][0][]0AIP}{[1][0][]source}{}{}{}}
\makeatother
\crefname{originallemma}{original lemma}{original lemmas}
\Crefname{originallemma}{Original lemma}{Original lemmas}
\makeatletter
\newlabel{lemma-presentation-representable}{{0AIG}{source}{}{sourcepointer-7}{}}
\newlabel{lemma-presentation-representable@cref}{{[originallemma][0][]0AIG}{[1][0][]source}{}{}{}}
\makeatother
\crefname{originallemma}{original lemma}{original lemmas}
\Crefname{originallemma}{Original lemma}{Original lemmas}
\makeatletter
\newlabel{lemma-separated-from-separated}{{0APX}{source}{}{sourcepointer-8}{}}
\newlabel{lemma-separated-from-separated@cref}{{[originallemma][0][]0APX}{[1][0][]source}{}{}{}}
\makeatother
\crefname{originallemma}{original lemma}{original lemmas}
\Crefname{originallemma}{Original lemma}{Original lemmas}
\makeatletter
\newlabel{more-groupoids-lemma-ind-quasi-affine}{{0APK}{source}{}{sourcepointer-9}{}}
\newlabel{more-groupoids-lemma-ind-quasi-affine@cref}{{[originallemma][0][]0APK}{[1][0][]source}{}{}{}}
\makeatother
\crefname{originallemma}{original lemma}{original lemmas}
\Crefname{originallemma}{Original lemma}{Original lemmas}
\makeatletter
\newlabel{more-morphisms-lemma-etale-separated-ind-quasi-affine}{{0AP9}{source}{}{sourcepointer-10}{}}
\newlabel{more-morphisms-lemma-etale-separated-ind-quasi-affine@cref}{{[originallemma][0][]0AP9}{[1][0][]source}{}{}{}}
\makeatother
\crefname{originallemma}{original lemma}{original lemmas}
\Crefname{originallemma}{Original lemma}{Original lemmas}
\makeatletter
\newlabel{morphisms-lemma-etale-locally-quasi-finite}{{03WS}{source}{}{sourcepointer-11}{}}
\newlabel{morphisms-lemma-etale-locally-quasi-finite@cref}{{[originallemma][0][]03WS}{[1][0][]source}{}{}{}}
\makeatother
\crefname{originallemma}{original lemma}{original lemmas}
\Crefname{originallemma}{Original lemma}{Original lemmas}
\makeatletter
\newlabel{spaces-properties-lemma-topology-points}{{03BX}{source}{}{sourcepointer-12}{}}
\newlabel{spaces-properties-lemma-topology-points@cref}{{[originallemma][0][]03BX}{[1][0][]source}{}{}{}}
\makeatother
\crefname{originalproposition}{original proposition}{original propositions}
\Crefname{originalproposition}{Original proposition}{Original propositions}
\makeatletter
\newlabel{spaces-properties-proposition-sheaf-fpqc}{{0APL}{source}{}{sourcepointer-13}{}}
\newlabel{spaces-properties-proposition-sheaf-fpqc@cref}{{[originalproposition][0][]0APL}{[1][0][]source}{}{}{}}
\makeatother
\crefname{originallemma}{original lemma}{original lemmas}
\Crefname{originallemma}{Original lemma}{Original lemmas}
\makeatletter
\newlabel{topologies-lemma-sheaf-property-fpqc}{{022H}{source}{}{sourcepointer-14}{}}
\newlabel{topologies-lemma-sheaf-property-fpqc@cref}{{[originallemma][0][]022H}{[1][0][]source}{}{}{}}
\makeatother
\crefname{originalremark}{original remark}{original remarks}
\Crefname{originalremark}{Original remark}{Original remarks}
\makeatletter
\newlabel{remark-set-theoretic}{{0AIS}{source}{}{sourcepointer-15}{}}
\newlabel{remark-set-theoretic@cref}{{[originalremark][0][]0AIS}{[1][0][]source}{}{}{}}
\makeatother
\title{A formal algebraic space satisfies the sheaf property for the fpqc topology}
\author{The Stacks Project Authors}
\date{stacks-project:formal-spaces.tex:local-snapshot-20261012; GNU FDL 1.2 or later}
\begin{document}
\maketitle
\noindent\textit{Editorial note.} All delivered mathematical text is the author's own, copied from the pinned Stacks Project chapter file formal-spaces.tex; no statement or proof step was rewritten, reordered or removed.\par\medskip
\noindent\textit{Editorial note.} The proof opens by saying that it is identical to the proof of a proposition in the chapter on properties of algebraic spaces. That sentence is an attribution, not a pointer in place of an argument: the proof is then displayed in full in the counted span, so the card is collected and the sentence is reproduced verbatim.\par\medskip
\noindent\textit{Editorial note.} The statement carries a slogan environment in this chapter; the chapter's own preamble defines slogan as a comment environment, so its contents are discarded at build time and are not part of the delivered mathematical text. The slogan is reproduced in the source span as delivered.\par\medskip
\noindent\textit{Editorial note.} The proof contains two steps the author marks as omitted: a small argument that the chosen etale cover may be taken to contain the image of the given morphism, and a computation that the two restrictions of the descended morphism agree. Both are reproduced verbatim and both are recorded as obligations of the card.\par\medskip
\noindent\textit{Editorial note.} The proof contains a parenthetical addressed to the reader, offering the alternative of replacing the chosen cover by a disjoint union of formal affines so that it is surjective; it lies inside the counted span and is reproduced verbatim.\par\medskip
\noindent\textit{Editorial note.} The counted spans carry 15 navigation stubs; no bibliography citation occurs in the spans.\par\medskip
\noindent\textit{Editorial note.} The preamble span points at the chapter's own preamble.tex shipped with the frozen source, used verbatim as staged.\par\medskip
\noindent\textit{Allowed standard prerequisites (editorial).}\begin{itemize}
\item The definition of a formal algebraic space, and the diagonal being a representable monomorphism. Same chapter, Definition definition-formal-algebraic-space at proof lines 7093-7095 and Lemma lemma-diagonal-formal-algebraic-space at proof lines 7076-7077. Fixes the objects and supplies the representable monomorphism that the uniqueness argument is built on.
\item Etale presentations: representability and separatedness of a morphism from a formal affine cover. Same chapter, Lemma lemma-presentation-representable at proof lines 7110-7111 and Lemma lemma-separated-from-separated at proof lines 7111-7112. Lets the base change of the chosen cover be treated as an etale separated morphism of schemes.
\item Etale separated morphisms are ind-quasi-affine, and ind-quasi-affine descent data are effective. More on Morphisms, Lemma more-morphisms-lemma-etale-separated-ind-quasi-affine at proof lines 7134-7136 (with Morphisms, Lemma morphisms-lemma-etale-locally-quasi-finite) and More on Groupoids, Lemma more-groupoids-lemma-ind-quasi-affine at proof lines 7136-7138. Supplies the effectivity of the descent datum built from the cocycle condition, which is what produces a scheme over the base.
\item Descent of surjectivity, of etaleness and of being an isomorphism, and fpqc coverings as universal effective epimorphisms. Descent, Lemma descent-lemma-descending-property-isomorphism at proof lines 7095-7097, and Descent, Lemmas descent-lemma-descending-property-surjective and descent-lemma-descending-property-etale at proof lines 7143-7145, and Descent, Lemma descent-lemma-fpqc-universal-effective-epimorphisms at proof lines 7153-7154. Carries the properties back along the covering and descends the morphism to the base.
\item The sheaf property for the fpqc topology is equivalent to the equalizer condition, and the corresponding statement for algebraic spaces. Topologies, Lemma topologies-lemma-sheaf-property-fpqc at proof lines 7066-7068, and Properties of Spaces, Proposition spaces-properties-proposition-sheaf-fpqc together with Properties of Spaces, Lemma spaces-properties-lemma-topology-points at proof lines 7061-7064 and 7101-7103. Fixes the reduction to the equalizer condition and supplies the model proof the author follows.
\end{itemize}\medskip
\ifdefined\hypertarget\hypertarget{sourcepointer-1}{}\fi
\noindent\textit{Source reference.} Original definition 023V: sources/stacks-chow/tags.txt, source lines 2756--2756. Pointer outside the counted span, invoked once for context; the load-bearing argument of the counted statement is the author's proof reproduced in full in the proof span.\par\medskip
\ifdefined\hypertarget\hypertarget{sourcepointer-2}{}\fi
\noindent\textit{Source reference.} Original lemma 02VN: sources/stacks-chow/tags.txt, source lines 3774--3774. Pointer outside the counted span, invoked once for context; the load-bearing argument of the counted statement is the author's proof reproduced in full in the proof span.\par\medskip
\ifdefined\hypertarget\hypertarget{sourcepointer-3}{}\fi
\noindent\textit{Source reference.} Original lemma 02L4: sources/stacks-chow/tags.txt, source lines 3395--3395. Pointer outside the counted span, invoked once for context; the load-bearing argument of the counted statement is the author's proof reproduced in full in the proof span.\par\medskip
\ifdefined\hypertarget\hypertarget{sourcepointer-4}{}\fi
\noindent\textit{Source reference.} Original lemma 02KV: sources/stacks-chow/tags.txt, source lines 3386--3386. Pointer outside the counted span, invoked once for context; the load-bearing argument of the counted statement is the author's proof reproduced in full in the proof span.\par\medskip
\ifdefined\hypertarget\hypertarget{sourcepointer-5}{}\fi
\noindent\textit{Source reference.} Original lemma 023Q: sources/stacks-chow/tags.txt, source lines 2751--2751. Pointer outside the counted span, invoked once for context; the load-bearing argument of the counted statement is the author's proof reproduced in full in the proof span.\par\medskip
\ifdefined\hypertarget\hypertarget{sourcepointer-6}{}\fi
\noindent\textit{Source reference.} Original lemma 0AIP: sources/stacks-formal-spaces/pinned-original-source.tex, source lines 2667--2667. Pointer outside the counted span, invoked once for context; the load-bearing argument of the counted statement is the author's proof reproduced in full in the proof span.\par\medskip
\ifdefined\hypertarget\hypertarget{sourcepointer-7}{}\fi
\noindent\textit{Source reference.} Original lemma 0AIG: sources/stacks-formal-spaces/pinned-original-source.tex, source lines 2381--2381. Pointer outside the counted span, invoked once for context; the load-bearing argument of the counted statement is the author's proof reproduced in full in the proof span.\par\medskip
\ifdefined\hypertarget\hypertarget{sourcepointer-8}{}\fi
\noindent\textit{Source reference.} Original lemma 0APX: sources/stacks-formal-spaces/pinned-original-source.tex, source lines 3744--3744. Pointer outside the counted span, invoked once for context; the load-bearing argument of the counted statement is the author's proof reproduced in full in the proof span.\par\medskip
\ifdefined\hypertarget\hypertarget{sourcepointer-9}{}\fi
\noindent\textit{Source reference.} Original lemma 0APK: sources/stacks-chow/tags.txt, source lines 13478--13478. Pointer outside the counted span, invoked once for context; the load-bearing argument of the counted statement is the author's proof reproduced in full in the proof span.\par\medskip
\ifdefined\hypertarget\hypertarget{sourcepointer-10}{}\fi
\noindent\textit{Source reference.} Original lemma 0AP9: sources/stacks-chow/tags.txt, source lines 13467--13467. Pointer outside the counted span, invoked once for context; the load-bearing argument of the counted statement is the author's proof reproduced in full in the proof span.\par\medskip
\ifdefined\hypertarget\hypertarget{sourcepointer-11}{}\fi
\noindent\textit{Source reference.} Original lemma 03WS: sources/stacks-chow/tags.txt, source lines 5123--5123. Pointer outside the counted span, invoked once for context; the load-bearing argument of the counted statement is the author's proof reproduced in full in the proof span.\par\medskip
\ifdefined\hypertarget\hypertarget{sourcepointer-12}{}\fi
\noindent\textit{Source reference.} Original lemma 03BX: sources/stacks-chow/tags.txt, source lines 4364--4364. Pointer outside the counted span, invoked once for context; the load-bearing argument of the counted statement is the author's proof reproduced in full in the proof span.\par\medskip
\ifdefined\hypertarget\hypertarget{sourcepointer-13}{}\fi
\noindent\textit{Source reference.} Original proposition 0APL: sources/stacks-chow/tags.txt, source lines 13479--13479. Pointer outside the counted span, invoked once for context; the load-bearing argument of the counted statement is the author's proof reproduced in full in the proof span.\par\medskip
\ifdefined\hypertarget\hypertarget{sourcepointer-14}{}\fi
\noindent\textit{Source reference.} Original lemma 022H: sources/stacks-chow/tags.txt, source lines 2706--2706. Pointer outside the counted span, invoked once for context; the load-bearing argument of the counted statement is the author's proof reproduced in full in the proof span.\par\medskip
\ifdefined\hypertarget\hypertarget{sourcepointer-15}{}\fi
\noindent\textit{Source reference.} Original remark 0AIS: sources/stacks-formal-spaces/pinned-original-source.tex, source lines 2769--2769. Pointer outside the counted span, invoked once for context; the load-bearing argument of the counted statement is the author's proof reproduced in full in the proof span.\par\medskip
\begin{definition}
\label{definition-formal-algebraic-space}
Let $S$ be a scheme. We say a sheaf $X$ on $(\Sch/S)_{fppf}$ is a
{\it formal algebraic space} if there exist a family of maps
$\{X_i \to X\}_{i \in I}$ of sheaves such that
\begin{enumerate}
\item $X_i$ is an affine formal algebraic space,
\item $X_i \to X$ is representable by algebraic spaces and \'etale,
\item $\coprod X_i \to X$ is surjective as a map of sheaves
\end{enumerate}
and $X$ satisfies a set theoretic condition
(see Remark \ref{remark-set-theoretic}). A
{\it morphism of formal algebraic spaces}
over $S$ is a map of sheaves.
\end{definition}

\begin{lemma}
\label{lemma-sheaf-fpqc}
\begin{slogan}
Formal algebraic spaces are fpqc sheaves
\end{slogan}
Let $S$ be a scheme. Let $X$ be a formal algebraic space over $S$. Then
$X$ satisfies the sheaf property for the fpqc topology.
\end{lemma}

\begin{proof}
The proof is {\bf identical} to the proof of
Properties of Spaces, Proposition
\ref{spaces-properties-proposition-sheaf-fpqc}.
Since $X$ is a sheaf for the Zariski topology it
suffices to show the following. Given a surjective
flat morphism of affines $f : T' \to T$ we have:
$X(T)$ is the equalizer of the two maps $X(T') \to X(T' \times_T T')$.
See Topologies, Lemma \ref{topologies-lemma-sheaf-property-fpqc}.

\medskip\noindent
Let $a, b : T \to X$ be two morphisms such that $a \circ f = b \circ f$.
We have to show $a = b$. Consider the fibre product
$$
E = X \times_{\Delta_{X/S}, X \times_S X, (a, b)} T.
$$
By Lemma \ref{lemma-diagonal-formal-algebraic-space}
the morphism $\Delta_{X/S}$ is a representable monomorphism. Hence
$E \to T$ is a monomorphism of schemes. Our assumption that
$a \circ f = b \circ f$ implies that $T' \to T$ factors (uniquely) through $E$.
Consider the commutative diagram
$$
\xymatrix{
T' \times_T E \ar[r] \ar[d] & E \ar[d] \\
T' \ar[r] \ar@/^5ex/[u] \ar[ru] & T
}
$$
Since the projection $T' \times_T E \to T'$ is a monomorphism
with a section we conclude it is an isomorphism. Hence we conclude that
$E \to T$ is an isomorphism by
Descent, Lemma \ref{descent-lemma-descending-property-isomorphism}.
This means $a = b$ as desired.

\medskip\noindent
Next, let $c : T' \to X$ be a morphism such that the two compositions
$T' \times_T T' \to T' \to X$ are the same. We have to find a morphism
$a : T \to X$ whose composition with $T' \to T$ is $c$. Choose a
formal affine scheme $U$ and an \'etale morphism $U \to X$ such that the image
of $|U| \to |X_{red}|$ contains the image of $|c| : |T'| \to |X_{red}|$.
This is possible by
Definition \ref{definition-formal-algebraic-space},
Properties of Spaces, Lemma \ref{spaces-properties-lemma-topology-points},
the fact that a finite union of formal affine algebraic spaces is a
formal affine algebraic space, and the fact that $|T'|$ is quasi-compact
(small argument omitted). The morphism $U \to X$ is representable
by schemes (Lemma \ref{lemma-presentation-representable}) and
separated (Lemma \ref{lemma-separated-from-separated}). Thus
$$
V = U \times_{X, c} T' \longrightarrow T'
$$
is an \'etale and separated morphism of schemes. It is also surjective
by our choice of $U \to X$ (if you do not want to argue this you can
replace $U$ by a disjoint union of formal affine algebraic spaces so that
$U \to X$ is surjective everything else still works as well). The fact that
$c \circ \text{pr}_0 = c \circ \text{pr}_1$ means that we obtain a
descent datum on $V/T'/T$
(Descent, Definition \ref{descent-definition-descent-datum})
because
\begin{align*}
V \times_{T'} (T' \times_T T')
& =
U \times_{X, c \circ \text{pr}_0} (T' \times_T T') \\
& =
(T' \times_T T') \times_{c \circ \text{pr}_1, X} U \\
& =
(T' \times_T T') \times_{T'} V
\end{align*}
The morphism $V \to T'$ is ind-quasi-affine by
More on Morphisms, Lemma
\ref{more-morphisms-lemma-etale-separated-ind-quasi-affine}
(because \'etale morphisms are locally quasi-finite, see
Morphisms, Lemma \ref{morphisms-lemma-etale-locally-quasi-finite}).
By More on Groupoids, Lemma \ref{more-groupoids-lemma-ind-quasi-affine}
the descent datum is effective. Say $W \to T$ is a morphism
such that there is an isomorphism $\alpha : T' \times_T W \to V$
compatible with the given descent datum on $V$ and the canonical descent
datum on $T' \times_T W$. Then $W \to T$ is surjective and \'etale
(Descent, Lemmas \ref{descent-lemma-descending-property-surjective} and
\ref{descent-lemma-descending-property-etale}).
Consider the composition
$$
b' : T' \times_T W \longrightarrow V = U \times_{X, c} T' \longrightarrow U
$$
The two compositions
$b' \circ (\text{pr}_0, 1), 
b' \circ (\text{pr}_1, 1) :
(T' \times_T T') \times_T W \to T' \times_T W \to U$
agree by our choice of $\alpha$ and the corresponding property of $c$
(computation omitted). Hence $b'$ descends to a morphism $b : W \to U$ by
Descent, Lemma \ref{descent-lemma-fpqc-universal-effective-epimorphisms}.
The diagram
$$
\xymatrix{
T' \times_T W \ar[r] \ar[d] & W \ar[r]_b & U \ar[d] \\
T' \ar[rr]^c &  & X
}
$$
is commutative. What this means is that we have proved the existence
of $a$ \'etale locally on $T$, i.e., we have an $a' : W \to X$.
However, since we have proved uniqueness
in the first paragraph, we find that this \'etale local solution
satisfies the glueing condition, i.e., we have
$\text{pr}_0^*a' = \text{pr}_1^*a'$ as elements of $X(W \times_T W)$.
Since $X$ is an \'etale sheaf we find a unique $a \in X(T)$ restricting
to $a'$ on $W$.
\end{proof}

\end{document}
